\documentclass[11pt,a4paper]{article}
\usepackage[utf8]{inputenc}
\usepackage[T1]{fontenc}
\usepackage[portuguese]{babel}
\usepackage{xcolor}
\usepackage{hyperref}
\usepackage{geometry}
\usepackage{listings}
\usepackage{tcolorbox}
\geometry{margin=2.5cm}
\definecolor{primary}{RGB}{41,128,185}
\definecolor{secondary}{RGB}{44,62,80}
\definecolor{codebg}{RGB}{240,240,240}
\lstset{backgroundcolor=\color{codebg},basicstyle=\ttfamily\small,breaklines=true,frame=single}
\title{\color{secondary}\Huge\textbf{StatefulSets, Volumes Persistentes e Jobs}}
\author{\color{primary}DevOps com Kubernetes -- 2026}
\date{}
\begin{document}
\maketitle


\section*{\color{primary}O que você aprenderá nesta página}
\begin{itemize}
\item Persistir dados com PV, PVC e StorageClass
\item Entender StatefulSets e Headless Services
\item Executar tarefas finitas com Jobs
\end{itemize}

Até aqui, os Pods podiam ser substituídos sem grande preocupação. Isso funciona bem para aplicações sem estado, mas não para bancos de dados, filas, caches persistentes e serviços que precisam de identidade estável. Nesses casos, precisamos combinar armazenamento persistente, \texttt{StatefulSet} e, em alguns cenários, \texttt{Job}.

\section*{\color{primary}Volumes persistentes}

\texttt{PersistentVolume} (PV) representa uma unidade de armazenamento disponível no cluster. \texttt{PersistentVolumeClaim} (PVC) é o pedido feito pela aplicação. \texttt{StorageClass} descreve como o armazenamento pode ser provisionado dinamicamente.

\begin{lstlisting}[language=bash]
apiVersion: v1
kind: PersistentVolumeClaim
metadata:
  name: app-data
spec:
  accessModes:
    - ReadWriteOnce
  resources:
    requests:
      storage: 1Gi
\end{lstlisting}

Um Pod ou Deployment pode montar esse PVC:

\begin{lstlisting}[language=bash]
volumeMounts:
  - name: data
    mountPath: /usr/share/nginx/html
volumes:
  - name: data
    persistentVolumeClaim:
      claimName: app-data
\end{lstlisting}

\section*{\color{primary}StatefulSet}

\texttt{StatefulSet} é o controlador indicado quando a aplicação precisa de identidade estável. Os Pods são criados com nomes previsíveis, como \texttt{postgres-0}, \texttt{postgres-1} e \texttt{postgres-2}. Um StatefulSet normalmente trabalha com um Headless Service para oferecer nomes DNS individuais para cada Pod.

\begin{lstlisting}[language=bash]
apiVersion: apps/v1
kind: StatefulSet
metadata:
  name: postgres
spec:
  serviceName: postgres-headless
  replicas: 1
  selector:
    matchLabels:
      app: postgres
  template:
    metadata:
      labels:
        app: postgres
    spec:
      containers:
        - name: postgres
          image: postgres:16
          ports:
            - containerPort: 5432
          volumeMounts:
            - name: data
              mountPath: /var/lib/postgresql/data
  volumeClaimTemplates:
    - metadata:
        name: data
      spec:
        accessModes: ["ReadWriteOnce"]
        resources:
          requests:
            storage: 1Gi
\end{lstlisting}

\section*{\color{primary}Jobs}

\texttt{Job} executa uma tarefa finita até a conclusão. Ele é útil para migrações, cargas iniciais, scripts administrativos e processamentos pontuais. Use Deployment para aplicações sem estado, StatefulSet quando identidade e volume por réplica forem importantes, e Job para tarefas que terminam.

\end{document}
